% Quelle des Templates: https://de.wikipedia.org/wiki/LaTeX

\documentclass[12pt]{exam}
\usepackage[T1]{fontenc}
\usepackage{lmodern}
\usepackage[ngerman]{babel}
\usepackage{amsmath}
\usepackage{graphicx}
\usepackage{tabularx}
\usepackage{tikz}
\usepackage{pgfplots}
\usepackage{amssymb}
\usepackage{trfsigns}
\usetikzlibrary{positioning}
\usepackage{hyperref}
\usepackage{verbatim}
\usepackage{listings}
\usepackage[siunitx,european]{circuitikz}
\usepackage{subcaption}
\usepackage{pythonhighlight}
\usepackage[utf8]{inputenc}

\newcommand{\emptybox}[2][\textwidth]{%
  \begingroup
  \setlength{\fboxsep}{-\fboxrule}%
  \noindent\framebox[#1]{\rule{0pt}{#2}}%
  \endgroup
}

\begin{document}

\begin {center}
Technische Hochschule\\
Würzburg-Schweinfurt \\
Fakultät Elektrotechnik\\
Sommersemester 2026 \\
\end{center}

\begin{tabular}{ll}
Sprachsteuerung:& Studiengang Elektro- und Informationstechnik (Bachelor) \\
Prozessdatenverarbeitung:& Studiengang Elektro- und Informationstechnik (Bachelor) \\
Speech Recognition:& Studiengang Robotik (Bachelor) \\
&\\
Datum: & 07.05.2026 \\
Dauer: & 90 Minuten \\
Hilfsmittel: & ein DIN A4 Blatt Formelsammlung (beidseitig) und Taschenrechner \\
Anzahl der Aufgaben: & 5\\
&\\
\end{tabular}

\begin{tabular}{ll}
Name:&...........................................\\
&\\
Matrikelnummer:&...........................................\\
&\\
Studiengang:&...........................................\\
&\\
Erstprüfer:&Prof. Dr.-Ing. Martin Spiertz\\
Zweitprüfer:&\\
&\\
\end{tabular}

\begin{tabular}{l}
Wichtige Hinweise: \\
Bearbeiten Sie die Aufgaben ausschließlich in den jeweiligen Lücken im Arbeitsblatt.\\
Sollte zusätzlicher Platz benötigt werden, werden Zusatzblätter ausgeteilt.\\
Kennzeichnen Sie auch alle Zusatzblätter mit Ihrem Namen.
\end {tabular}

\begin{questions}
\clearpage %Aufgabe 1
\question Aufgabe (10 Punkte):\\
\begin{parts}
\part Geben Sie die mathematische Beschreibung eines digitalen Signals an, dass den maximalen Pegel in dB FS hat. Berechnen Sie den dazugehörigen Pegel in dB FS?\\\\
\begin{centering}
\begin{tikzpicture} %Bild 3
\draw[step =.5cm ,gray ,very  thin] (0,0)  grid(15 ,5);
\end{tikzpicture}
\end{centering}\\\\
Ein Signal $x(t)=\cos\left(500\cdot t\right)$ wird mit $r=249$ Hz abgetastet.
\part Verursacht diese Abtastung Alias?\\\\
\begin{centering}
\begin{tikzpicture} %Bild 3
\draw[step =.5cm ,gray ,very  thin] (0,0)  grid(15 ,2);
\end{tikzpicture}
\end{centering}\\\\
\part Berechnen Sie das abgetastete Signal für $n=2$. Hinweis: $x(n)=x\left(t=n/r\right)$\\\\
\begin{centering}
\begin{tikzpicture} %Bild 3
\draw[step =.5cm ,gray ,very  thin] (0,0)  grid(15 ,7);
\end{tikzpicture}
\end{centering}\\\\
\clearpage
Ein FIR Filter ist gegeben: $h(n)=\left[1; -1\right]$
\part Ist $h(n)$ stabil? (Begründen Sie Ihre Antwort)\\\\
\begin{centering}
\begin{tikzpicture} %Bild 3
\draw[step =.5cm ,gray ,very  thin] (0,0)  grid(15 ,2);
\end{tikzpicture}
\end{centering}\\\\
\part Berechnen Sie den Realteil der Übertragungsfunktion $H(f)$ von $h(n)$ mit Hilfe der z-Transformation. Ist der Realteil für ein beliebiges $f$ negativ?\\\\
\begin{centering}
\begin{tikzpicture} %Bild 3
\draw[step =.5cm ,gray ,very  thin] (0,0)  grid(15 ,5);
\end{tikzpicture}
\end{centering}\\\\
Ein Signal $y(n)$ wird durch folgenden Python Code erzeugt:
\begin{python}
import numpy as np
y = np.random.randn(10000)
\end{python}
\part Bestimmen Sie den Python Code um den Mittelwert und die Standardabweichung eines Blocks von $y(n)$ beginnend bei $n=1000$ und eine Blocklänge von $N=1000$ zu bestimmen.\\\\
\begin{centering}
\begin{tikzpicture} %Bild 3
\draw[step =.5cm ,gray ,very  thin] (0,0)  grid(15 ,5);
\end{tikzpicture}
\end{centering}\\\\
\end{parts}

\clearpage %Aufgabe 2
\question Aufgabe (10 Punkte):\\
\begin{parts}
\part Geben Sie zwei Gründe für die Benutzung von Zeropadding bei einer DFT an.\\\\
\begin{centering}
\begin{tikzpicture} %Bild 3
\draw[step =.5cm ,gray ,very  thin] (0,0)  grid(15 ,5);
\end{tikzpicture}
\end{centering}\\\\
\part Warum sollte das Eingangssignal der DFT mit einer Fensterfunktion multipliziert werden bevor man die DFT anwendet?\\\\
\begin{centering}
\begin{tikzpicture} %Bild 3
\draw[step =.5cm ,gray ,very  thin] (0,0)  grid(15 ,5);
\end{tikzpicture}
\end{centering}\\\\
\clearpage
Im Folgenden wird eine STFT auf ein Signal $x(n)$ angewendet. Die Abtastrate ist $r=48$ kHz, das Signal hat eine Dauer von $70$ Minuten. Die Blocklänge ist $N=1000$, die Überlappung ist $75$ \%, die Transformationslänge ist $K=2028$.
\part Was sind die Vorteile von großen Blocklängen $N$? Was sind die Vorteile von kurzen Blocklängen $N$?\\\\
\begin{centering}
\begin{tikzpicture} %Bild 3
\draw[step =.5cm ,gray ,very  thin] (0,0)  grid(15 ,5);
\end{tikzpicture}
\end{centering}\\\\
\part Bestimmen Sie die Größe der Matrix des resultierenden Spektrogramms.\\\\
\begin{centering}
\begin{tikzpicture} %Bild 3
\draw[step =.5cm ,gray ,very  thin] (0,0)  grid(15 ,7);
\end{tikzpicture}
\end{centering}\\\\
\part Was bedeuten die Abkürzungen DFT, FFT und STFT?\\\\
\begin{centering}
\begin{tikzpicture} %Bild 3
\draw[step =.5cm ,gray ,very  thin] (0,0)  grid(15 ,2);
\end{tikzpicture}
\end{centering}
\end{parts}

\clearpage %Aufgabe 3
\question Aufgabe (10 Punkte):\\
\begin{parts}
\part Ein Signal hat die Frequenz $f=1234$ Hz. Bestimmen Sie die dazugehörigen Werte in Bark und mel.\\\\
\begin{centering}
\begin{tikzpicture} %Bild 3
\draw[step =.5cm ,gray ,very  thin] (0,0)  grid(15 ,5);
\end{tikzpicture}
\end{centering}\\\\
\part Was ist die Zeitauflösung in Millisekunden und die Frequenzauflösung in Bark für das menschliche Gehör?\\\\
\begin{centering}
\begin{tikzpicture} %Bild 3
\draw[step =.5cm ,gray ,very  thin] (0,0)  grid(15 ,5);
\end{tikzpicture}
\end{centering}\\\\
\clearpage
Im Folgenden gilt: $y(n)=x(n)+d(n)$, $x(n)$ und $d(n)$ sind unkorreliert. Bei $d(n)$ handelt es sich um additives weißes Rauschen.
\part Für welches $p$ ist folgende Gleichung erfüllt: $\left|Y(f)\right|^p\approx\left|X(f)\right|^p+\left|D(f)\right|^p$\\\\
\begin{centering}
\begin{tikzpicture} %Bild 3
\draw[step =.5cm ,gray ,very  thin] (0,0)  grid(15 ,2);
\end{tikzpicture}
\end{centering}\\\\
\part Wie lautet die Abkürzung für das gegebene Rauschmodell?\\\\
\begin{centering}
\begin{tikzpicture} %Bild 3
\draw[step =.5cm ,gray ,very  thin] (0,0)  grid(15 ,3);
\end{tikzpicture}
\end{centering}\\\\
\part Welche Eigenschaften erfüllt das Spektrum von $d(n)$?\\\\
\begin{centering}
\begin{tikzpicture} %Bild 3
\draw[step =.5cm ,gray ,very  thin] (0,0)  grid(15 ,3);
\end{tikzpicture}
\end{centering}\\\\
\part Welche Eigenschaften erfüllen die Abtastwerte $d(n)$ im Zeitbereich?\\\\
\begin{centering}
\begin{tikzpicture} %Bild 3
\draw[step =.5cm ,gray ,very  thin] (0,0)  grid(15 ,5);
\end{tikzpicture}
\end{centering}
\end{parts}

\clearpage %Aufgabe 4
\question Aufgabe (10 Punkte):\\
\begin{parts}
\part Menschliche Sprache wird mit einem Mikrofon aufgenommen. Was ist der typische Frequenzbereich der menschlichen Sprache? Was ist der typische Mittelwert des aufgenommenen Signals?\\\\
\begin{centering}
\begin{tikzpicture} %Bild 3
\draw[step =.5cm ,gray ,very  thin] (0,0)  grid(15 ,6);
\end{tikzpicture}
\end{centering}\\\\
\part Welches Delay erzeugt ein symmetrisches FIR-Filter im Durchlassbereich in Millisekunden? Die Länge des Filters beträgt $N=20$, die Abtastrate beträgt $r=16$ kHz.\\\\
\begin{centering}
\begin{tikzpicture} %Bild 3
\draw[step =.5cm ,gray ,very  thin] (0,0)  grid(15 ,6);
\end{tikzpicture}
\end{centering}\\\\
\clearpage
Im Folgenden wird das IIR-Filter $y(n)=a\cdot y(n-1)+(1-a)\cdot x(n)$ betrachtet.
\part Für welche Werte von $a$ ist das Filter stabil?\\\\
\begin{centering}
\begin{tikzpicture} %Bild 3
\draw[step =.5cm ,gray ,very  thin] (0,0)  grid(15 ,3);
\end{tikzpicture}
\end{centering}\\\\
\part Geben Sie die Verstärkung des Gleichanteils an: $H\left(f=0\right)$.\\\\
\begin{centering}
\begin{tikzpicture} %Bild 3
\draw[step =.5cm ,gray ,very  thin] (0,0)  grid(15 ,5);
\end{tikzpicture}
\end{centering}\\\\
\part Ein Eingangssignal $x(n)$ ist gegeben. Ergänzen Sie den Python Code, um das obige IIR-Filter zu implementieren.\\\\
\begin{python}
import numpy as np
a = 0.9
N = 1000
x = np.random.randn(N)
\end{python}
\begin{centering}
\begin{tikzpicture} %Bild 3
\draw[step =.5cm ,gray ,very  thin] (0,0)  grid(15 ,5);
\end{tikzpicture}
\end{centering}
\end{parts}

\clearpage %Aufgabe 5
\question Aufgabe (10 Punkte):\\
Die folgenden zehn Python-Schnipsel sollen alle mit einem vorangestellten
\begin{python}
import numpy as np
\end{python}
beginnen. Bestimmen Sie für jeden Schnipsel, ob die Ausgabe A oder B ist (durch Markieren der korrekten Ausgabe auf dem Aufgabenblatt).

\begin{python}
x = np.ones(5)
assert x.sum()==5, `A'
print(`B')
\end{python}

\begin{python}
x = np.array([1, 0, 0])
a = np.array([1, -0.5])
y = np.array([0, 0])
for i in range(2):
    y[i] = x[i] + 0.5 * y[i-1] if i > 0 else x[i]
assert (y == [1, 0.5]).all(), `A'
print(`B')
\end{python}

\begin{python}
x = np.array([1, 0, 0])
b = np.array([1, 1, 1])
y = np.convolve(x, b, mode='full')
assert (y == [1, 1, 1, 0, 0]).all(), `A'
print(`B')
\end{python}

\begin{python}
x = np.array([1, 0, 0, 0])
X = np.fft.fft(x)
assert X[0] == X[-1], `A'
print(`B')
\end{python}

\begin{python}
x = np.array([1, 2, 3])
b = np.array([1, 1])
y = np.convolve(x, b, mode='valid')
assert (y == [3, 5]).all(), `A'
print(`B')
\end{python}

\begin{python}
x = np.array([1, 2, 3, 4])
X = np.fft.fft(x)
assert np.isreal(X).all(), `A'
print(`B')
\end{python}

\begin{python}
x = np.array([1, 1, 1, 1])
X = np.fft.fft(x)
assert abs(X[1]) == 0, `A'
print(`B')
\end{python}

\begin{python}
x = np.array([1, 2, 3])
assert np.var(x) == 0.6666666666666666, `A'
print(`B')
\end{python}

\begin{python}
x = np.array([1, 2, 3])
assert np.std(x) == 1, `A'
print(`B')
\end{python}

\begin{python}
x = np.array([1, 2, 3])
assert np.mean(x) == 2, `A'
print(`B')
\end{python}

\end{questions}
\end{document}



